\section{与 EP139、EP136 的连接}

EP139 讲 Agent 技术史，强调 proactive agent 和 universal digital agent；EP136 讲模型 OS，强调模型成为任务调度层；EP135 则把这两个判断放到社交产品里：如果 AI 能主动完成任务，那么它也能主动帮人建立关系。如果模型是 OS，那么社交网络可能成为 OS 上最敏感的个人 context 应用。

\begin{knowledgebox}{三期连读框架}
EP139 给出 Agent 的技术定义，EP136 给出模型 OS 的产业图景，EP135 给出 AI 社交的产品实例。它们共同回答：模型如何从回答问题，走向代表人行动。
\end{knowledgebox}

\subsection{本章小结}

EP135 是前两期的 C 端产品化延伸。它不是讨论模型能力本身，而是讨论当模型可以主动行动时，社交网络会怎样重写。

